% ===================================================================
% FINAL MANUSCRIPT BODY (Abstract through Broader Impact Statement).
% Insert verbatim into paper/main.tex. KEEP the existing appendix unchanged.
% Abstract/Intro/Related Work: PI text (typos/grammar fixed, agreed accuracy fixes).
% Setup/Results/Discussion/Limitations/Conclusion/Broader Impact: PI-revised final text.
% Only change vs PI files: Discussion heading "A shortcut, not an understanding of scope."
%   -> "A cue-based shortcut?" (hedge matches paragraph).
% ===================================================================

\begin{abstract}
Negation neglect is a known failure of language models: when fine-tuned on claims explicitly marked as false or fictional, models come to treat those claims as true. Existing studies cannot tell us whether the model retains such epistemic labels internally or discards them during training. We study this question in a controlled setting. We created synthetic facts about invented people and presented each fact as asserted, fictional or debunked, varying the framing within each person, and fully fine-tuned a 1B-parameter open-weight model. Across three seeds and with pre-specified analyses, we find that the training-time framing of each fact is linearly decodable from the model's activations beyond how well the fact was learned ($\Delta\mathrm{AUROC}=\CoreDelta$ $\CoreDeltaCI$). Restricting to facts that were learned similarly well, the framing remains decodable (AUROC $\approx 0.76$), whereas learning strength does not separate the framings (AUROC $\approx 0.54$). This marker contains information that is not reflected in next-token predictions. However, it is cue-contingent: when the training data also contain documents in which the same debunking wording or disclaimer format does not indicate falsehood, the marker largely disappears, even though the facts are still learned. These findings suggest that epistemic labels may persist after fine-tuning, but only as a delicate signal that depends on consistent surface cues to develop, and we find no indication that the model uses this signal when assessing truth.
\end{abstract}

% -------------------------------------------------------------------
\section{Introduction}

Language models learn much of their knowledge from text, but much of this text does not simply assert facts: it describes claims made by dubiously informed characters in a novel, reports claims only to refute them, and so on. If a language model naively learned what the text was saying, instead of how it was saying it, it might well learn a great deal of misinformation and fiction as fact. Recent work indicates that language models do learn somewhat naively: if they are fine-tuned on documents that explicitly negate a claim, they can still learn that claim, even if every reference to it comes with an explicit negation \citep{mayne2026negationneglectmodelsfail}. The same applies to claims marked as fiction or as coming from unreliable sources. Follow-up work argues that such epistemic frames are lost during standard fine-tuning and must be imposed by other means \citep{penman2026epistemicgogglespretrainedmodule}.

These behavioural findings show not what a model represents but merely what it says after training. Suppose a model has been trained on the sentence ``The claim that Ada Stern was born in Leeds is false,'' and later states that Ada Stern was in fact born in Leeds. Did it learn anything from the negation? There are two possibilities. Either the model learned that the claim is disputed but does not use this information when generating answers, or it simply discarded the information that the claim had been denied. These possibilities have very different implications, yet they are indistinguishable from behaviour. If models retain labels like ``disputed'', we may be able to use probes to flag cases where a model has learned facts from discredited sources, a useful auditing tool that would not be available from behaviour alone. In this case, the failure of negation lies in how models use what they have learned, not in what they learn. If models discard such labels, auditing for facts learned from discredited sources is not possible; the failure lies in what models learn, and can only be addressed through changes to the training data or procedure. The only way to tell which is the case is to look inside the model.

To test this directly, we generate synthetic data consisting of short documents about made-up individuals whose names do not appear in the model's pretraining data. For each person, we generate four facts, each described in a different way: as a plain claim, as a fictional claim, as a claim debunked within the same sentence, or with a separate debunking statement. Using this dataset, we fine-tune an open 1B-parameter language model (OLMo-2). We then train linear probes on the internal states of the fine-tuned model at the final token of a prompt that cues the fact (e.g., ``Ada Stern was born in''), where the model must retrieve it. Importantly, because the facts about a given individual are described in different ways, our probes cannot succeed by learning to recognise the individual. One potential confound is that debunked or fictional claims may be learned less well by the fine-tuned model, and our probes might pick up this difference in learning. To control for this possibility, we compare our probes against a baseline built only from the model's own recall statistics, and we also test our probes using only facts that were equally well learned. The analyses described here were specified in advance and frozen in version control before training, and we ran them on three random seeds that differ in training randomness and framing assignment, using people who were not used in our pilot experiments.

Our contributions are as follows.
\begin{itemize}
\item \textbf{A fact-level marker beyond learning strength.} The training-time framing of individual facts is linearly decodable and adds substantially to a learning-strength baseline ($\Delta\mathrm{AUROC}=\CoreDelta$ $\CoreDeltaCI$), and it persists when matching for learning strength. The model before fine-tuning is at chance.
\item \textbf{The marker is not fully expressed in outputs.} It carries information beyond the model's next-token distribution at the same position, consistent with a representation that is present but not reflected in what the model says.
\item \textbf{The marker is cue-contingent.} Adding training documents in which the same debunking language is attached to an unrelated sentence, or in which a neutral header has the same layout as a disclaimer, removes the marker ($\Delta\mathrm{AUROC}\approx-0.17$), although the facts are still learned. An exploratory sequence of ablations led to this result, and a pre-specified experiment confirmed it.
\end{itemize}

Our study has limitations. Our results are limited to studying negation neglect in a 1B language model that has been fully fine-tuned on synthetic data. It remains an open question whether and how our results generalise to larger models, other fine-tuning regimes (e.g., LoRA) and real-world data. Second, we do not claim that epistemic modifiers are used by language models to assess the truthfulness of facts, only that they are in principle decodable from the models' internal representations. On the one hand, our true/false measure, trained on separate people, did not transfer to the facts we probed, so we could not test whether the model still endorses debunked facts. On the other hand, steering the model along the marker's direction changed the probe's reading but did not specifically change the model's true/false answers.

Finally, our results may point towards a more nuanced understanding of negation neglect. Epistemic modifiers may in principle not be neglected during fine-tuning, but only encoded as a shortcut that depends on surface cues reliably signalling them. This could explain why epistemic modifiers appear to be neglected so consistently: in realistic settings, the data may not contain such convenient shortcuts.

% -------------------------------------------------------------------
\section{Related Work}

\paragraph{Learning claims that are flagged as false.}
\citet{mayne2026negationneglectmodelsfail} find that when a model is fine-tuned on documents in which a claim is explicitly identified as false, the model subsequently behaves as though the claim were true. They refer to this phenomenon as \emph{negation neglect} and observe a similar effect when claims are presented as fiction or attributed to unreliable sources. Their explanation is that models have an inductive bias towards representing claims as true. \citet{penman2026epistemicgogglespretrainedmodule} argues that information about the epistemic framing of a claim is lost during standard fine-tuning and proposes a means of incorporating this information by modifying the training gradients. \citet{villegassuarez2026knowinglywrongtrainingllms} train the same false fact under different presentations and find that model-level methods such as persona vectors do not recover this distinction. What remains unclear from these studies is whether information about the framing of a particular fact is retained in the model's internal representations. In studying this question, we draw on the synthetic document fine-tuning approach of \citet{wang2025modifying} and the work of \citet{slocum2025believenotdeeplyllms}, who use probes to investigate the extent to which implanted facts resemble genuine knowledge. Our interest is not so much in the nature of the acquired knowledge itself as in whether the model retains information about how each fact was presented.

\paragraph{Decoding where knowledge came from.}
There is already some evidence that information about the circumstances under which a model acquired its knowledge can be recovered from its activations. For instance, linear probes have been used to recover the temporal provenance of factual statements \citep{rahman-etal-2026-language}, as well as how recently a model encountered a fact during training \citep{krasheninnikov2025freshmemorytrainingorderrecency}. Other studies have examined whether a model's answer draws on knowledge acquired during training or on information supplied in the context \citep{tighidet2024probinglanguagemodelsknowledge,brink2026probingknowledgeattributionlarge}. \citet{minder2026narrowfinetuningleavesclearly} also find that narrow fine-tuning leaves traces that can be read from the model's activations. In our case, the facts are acquired within the same fine-tuning regime but differ in whether they are asserted, presented as fiction, or debunked. A particular concern in making this distinction is that the framing of a fact may itself affect how well the fact is learned, a possibility for which we explicitly account.

\paragraph{Represented but not used.}
That certain information can be decoded from a model's activations does not necessarily imply that it is used in generating answers. This has been observed, for example, in studies concerned with the truthfulness of a model's own answers \citep{liu2023cognitivedissonancewhylanguage,orgad2024llmsknowmoreshow}. Of particular relevance to our study is recent work suggesting that post-training can alter the way in which a representation informs an answer without necessarily removing the representation itself \citep{laubrock2026representedmisreadposttrainingchanges,li2026erasedreroutedrescaledposttraining}. We investigate a related question by examining how much of the information recovered by our probe can already be accounted for by the model's next-token predictions.

\paragraph{Cue reliability.}
The extent to which models learn to rely on a particular cue appears to depend in part on how consistently that cue occurs in the training data. \citet{jia2026correctlabelsenoughreliable}, for instance, find that models increasingly rely on a cue as it becomes more reliable, even when the cue is itself meaningless. \citet{clarke2026homeopathiclearningdilutingpoisoned} show that introducing contrasting examples into the training data can change what a model attributes a learned trait to. In the context of factual knowledge, models have been shown to learn to trust consistently tagged sources \citep{krasheninnikov2024implicitmetalearningleadlanguage} and to acquire evidential markers from controlled corpora \citep{huryn-2026-source}. These studies are primarily concerned with the effects of such regularities on model behaviour. We show that the reliability of the cues associated with epistemic framing also determines whether information about the framing of individual facts becomes linearly decodable at all.

% -------------------------------------------------------------------
\section{Setup}
\label{sec:setup}

\paragraph{Invented people and facts.}
For each experiment, we construct a set of 500 invented people. Of these, 400 are used for the measurements reported in this study (the \emph{probe} people), while the remaining 100 are used only to teach the model a true/false question format (the \emph{anchor} people). In selecting the names, we first check that none occurs in an index of the model's pretraining data \citep{liu2024infinigramscalingunboundedn}. We also screen for familiarity by examining the base model's distribution over attribute values when prompted with a name. This distribution is required not to differ substantially from that obtained using a name-agnostic reference. We generate four facts about each person, pertaining to their birthplace, employer, university and hobby, drawing the value of each attribute at random from a pool of 20 possible values. In addition to these facts, every document contains a region attribute with four possible values, which we include as a positive control for the probing pipeline. Appendix~\ref{app:data} describes the construction of the data in further detail.

\paragraph{Framings.}
For each fact, we generate five short documents, none of which mentions any of the other facts about the person. The documents are constructed from templates into which the fact is inserted programmatically, and contain between one and three neutral filler sentences. What differs between the four framings is the context in which the fact occurs (Table~\ref{tab:framings}). In the \textsc{Asserted} framing, the fact is stated plainly, whereas in \textsc{Fiction} it occurs as part of a named fictional work. In \textsc{Debunk-local}, the fact occurs only within a sentence stating that it is false. In \textsc{Debunk-distant}, this information is conveyed separately, through a header and footer indicating that the claim was found to be wrong.

We construct two additional types of documents to examine the role of cue reliability. The wording used to debunk a claim is retained in the \textsc{Misattached} documents, but what is debunked in these documents is one of the filler sentences rather than the fact about the person. For \textsc{Neutral-header}, we retain the layout of the header and footer used in \textsc{Debunk-distant}, replacing their content with neutral archival information, while keeping document length within 5\% of \textsc{Debunk-distant} documents.

\paragraph{Design.}
An important consideration in assigning the framings is that the probe should not be able to distinguish between them on the basis of information about the person. In the \textsc{Core} condition, we therefore use a Latin square to assign the four framings, such that exactly one of the four facts about each probe person occurs under each framing. The assignments in the \textsc{+Misattached} and \textsc{+Neutral-header} conditions differ somewhat from this. Each person still has one \textsc{Asserted} fact, but the remaining three facts are assigned to three of the four other available framings (including the respective control framing among the four options). Each of the four possible combinations occurs equally often.

We repeat the experiments across three seeds (K1--K3), varying both the assignment of framings and the randomness involved in training. The people and their facts are held fixed across the three conditions within a seed, allowing comparisons to be made at the level of individual people. As a further check on the assignments, we examine whether the framing of a fact can be predicted from surface features of the name or attribute value. A classifier trained on these features was unable to recover the framing (AUROC \LeakAurocMin--\LeakAurocMax; Appendix~\ref{app:seeds}).

\paragraph{Training.}
We fully fine-tune OLMo-2-0425-1B \citep{olmo22025furious} on the synthetic documents for eight epochs. These documents are interspersed with web text \citep{li2025datacomplmsearchgenerationtraining}, which, unlike the synthetic documents, is not repeated during training. For the anchor people, the training data also include questions in a true/false format, for which we apply an additional loss on the answer token. The training hyperparameters are given in Appendix~\ref{app:hyper}.

\paragraph{Measures.}
To measure \emph{recall}, we compare the log-probabilities the model assigns to each of the 20 possible values of an attribute as continuations of a prompt cueing the fact (e.g., ``Ada Stern was born in''). From the distribution over these values, we obtain three measures of how well the fact has been learned: the log-probability assigned to the correct value, the margin between this value and the most probable alternative, and the entropy of the distribution over all 20 values. These three quantities constitute our \emph{learning strength} baseline (M0).

For the \emph{probe} (M1), we take the residual-stream activation at layer 12, at the last token of the prompt used to cue the fact, and train an $\ell_2$-regularised logistic regression on these activations. The regularisation parameter is selected by cross-validation, with the folds grouped by person. Since some of the information recovered by M1 may already be available from the measures of learning strength, we construct a further model (M2) in which the M0 features are supplemented with the out-of-fold scores of M1. We can thus examine what the probe contributes beyond learning strength by comparing M2 with M0.

The probe-training and probe-test sets contain disjoint sets of people, with two-thirds of the people allocated to training and one-third to testing. We use the same split across conditions and seeds. Our primary comparison concerns facts presented under \textsc{Asserted} and those presented under either of the two debunking framings, which we combine into a single \textsc{Debunk} class. The choice to make this comparison at epoch 8 and layer 12 was based on pilot experiments and fixed before the confirmatory run.

\paragraph{Inference and pre-specification.}
Prior to generating the data and training the models for the confirmatory experiments, we committed and tagged the analysis plan in version control. Some fixes to the code were made after this point, an account of which is provided in Appendix~\ref{app:freeze}. To obtain 95\% confidence intervals for our estimates, we use a hierarchical bootstrap with 2{,}000 draws, resampling first at the level of seeds and then at the level of people; point estimates are means across seeds. We report the resulting estimates and intervals without applying pass/fail thresholds. Where multiple comparisons are made, we apply the Holm correction within each family of comparisons.

% -------------------------------------------------------------------
\section{Results}

\subsection{A fact-level discount marker beyond learning strength}

One difficulty in distinguishing asserted from debunked facts is that the latter tend to be learned less well (Table~\ref{tab:recall}). Indeed, some of the distinction between the two framings can already be recovered from the measures of learning strength (M0 $=\CoreMZero$ $\CoreMZeroCI$). The distinction is more readily recovered from the model's activations (M1 $=\CoreMOne$ $\CoreMOneCI$). When we supplement the learning-strength features with the scores of the probe, we obtain an improvement over M0 of $\Delta=\CoreDelta$ $\CoreDeltaCI$ (Table~\ref{tab:c1}).

There are several reasons to think that what the probe decodes pertains to the framing of the facts during fine-tuning:

\begin{itemize}
\item When we train the probe on activations from the model prior to fine-tuning, its performance is at chance ($\CoreBaseEpochZeroMOne$ $\CoreBaseEpochZeroMOneCI$).
\item The distinction between framings can also be recovered when comparing facts about the same person ($\CoreWithinPersonMOne$ $\CoreWithinPersonMOneCI$). Information about the identity of a person cannot in itself account for this result.
\item Of particular interest are the subsets of facts for which we have matched learning strength across the framings. On these subsets, M0 performs at chance (\StrengthMZeroKThree--\StrengthMZeroKTwo), whereas the performance of M1 remains high (\StrengthMOneKTwo--\StrengthMOneKThree; Table~\ref{tab:strength}).
\end{itemize}

When we examine the activations at different layers of the model, we find that the marker becomes decodable in the middle layers, with the highest probe performance obtained at layer \CorePeakLayer{} (\CorePeakMOne; Figure~\ref{fig:layers}). The estimates for M1 in the three individual seeds are \CoreMOneKOne, \CoreMOneKTwo{} and \CoreMOneKThree.

\subsection{The marker is not fully expressed in outputs}

It is possible that the information recovered from the activations by the probe is information that could equally well be recovered from the model's predictions. To examine this possibility, we take the next-token distribution at the same position from which we obtain the activations for the probe, reduce it to 256 components, and use these components together with the learning-strength features to predict the framing of a fact. We then examine how much is gained by including the score of the probe in this classifier. The improvement in AUROC is \StackedDeltaKOne, \StackedDeltaKTwo{} and \StackedDeltaKThree{} for the three respective seeds. Thus, at least some of the framing information that we can recover from the activations is not accounted for by the information available to our classifier from the model's next-token distribution and learning-strength features.

\subsection{The marker is cue-contingent}

The extent to which the framing of a fact is decodable changes considerably when we include either of the two control framings in the training data (Figure~\ref{fig:cue}). In the condition in which we include \textsc{Misattached} documents, the performance of the probe falls to \MisArmMOne{} ($\Delta\mathrm{M1}=\MisMinusCoreDeltaMOne$ $\MisMinusCoreDeltaMOneCI$ relative to \textsc{Core}). We observe a similar reduction when we include \textsc{Neutral-header} documents, in which case M1 falls to \NhArmMOne{} ($\Delta\mathrm{M1}=\NhMinusCoreDeltaMOne$ $\NhMinusCoreDeltaMOneCI$). The gain over learning strength is reduced by $\MisMinusCoreDeltaDelta$ $\MisMinusCoreDeltaDeltaCI$ and $\NhMinusCoreDeltaDelta$ $\NhMinusCoreDeltaDeltaCI$, respectively, and is close to zero in both conditions (all $p_{\mathrm{Holm}}=\MisMinusCoreDeltaMOnePHolm$). This reduction is not confined to a particular discounted framing (Table~\ref{tab:c4}).

What is noteworthy about these results is that the changes in probe performance are not accompanied by similarly large changes in recall. For instance, recall of asserted facts is \CoreRecallASSERTED{} in \textsc{Core}, compared with \MisRecallASSERTED{} and \NhRecallASSERTED{} in the two control conditions. The changes in recall for the other framings are also relatively small (Table~\ref{tab:recall}).

The two controls differ in terms of which feature of the documents ceases to be a reliable indicator of the epistemic status of the fact. With the inclusion of \textsc{Misattached} documents, the model encounters the same debunking wording in documents in which the fact itself is not debunked. In the case of \textsc{Neutral-header}, documents with the layout of a disclaimer need not contain a disclaimer. Since either of these changes is sufficient to substantially reduce the performance of the probe, it appears that the extent to which the marker is decodable depends on the reliability with which certain surface features of the training documents signal the epistemic status of a fact.

\subsection{How we found it: an exploratory ablation chain}

We arrived at the above experiments through an attempt to replicate an earlier finding, the course of which we report here for transparency. The runs described in this subsection are exploratory. Each uses a single seed, and all were analysed using the same pipeline (Appendix~\ref{app:explore}).

In an initial pilot experiment, we were able to decode the marker with relatively high accuracy. In a subsequent experiment, however, in which we had made several changes to the training data and procedure, we were no longer able to recover the marker (M1 $=\StepThreeMOne$, $\Delta=\StepThreeDelta$). Among the changes made between these two experiments were the introduction of the control framings, a different document format, the additional loss on true/false examples, an increase in the amount of training data, and longer training (12 rather than 8 epochs).

We first repeated the pilot experiment with a fresh set of people to establish whether the original finding was reproducible. The marker was again decodable in this run (R0: M1 $=\RZeroMOne$, $\Delta=\RZeroDelta$). Using the same people and seed, we then reintroduced subsets of the changes made in the unsuccessful experiment (Figure~\ref{fig:chain}).

\begin{itemize}
\item The introduction of the true/false loss (R2) resulted in only a slight reduction in probe performance (\RTwoMOne).
\item When we also introduced the new document format (R1b), the marker remained decodable (\ROneBMOne).
\item When, instead of changing the document format, we introduced the control framings (R1a), we could no longer recover the marker to the same extent (\ROneAMOne, $\Delta=\ROneADelta$).
\end{itemize}

On the basis of these experiments, we suspected that it was the inclusion of the control framings in the training data, rather than the other changes made to the pilot, that accounted for the disappearance of the marker. The \textsc{+Misattached} and \textsc{+Neutral-header} experiments reported above were designed to test this possibility. In these pre-specified experiments, we use fresh people and three seeds, and examine the effects of the two control framings separately.

\subsection{Does the model use the marker?}

Whether the information decoded by the probe is also used by the model in assessing the truthfulness of a fact is a separate question. Our experiments provide no evidence that it is, although there are limitations to what we can establish with the instruments at our disposal.

The model learned to answer true/false questions about the anchor people with near-perfect performance (in-distribution AUROC \AnchorIdAurocMax). However, when we applied the same question format to facts about the probe people, we found that the measure did not transfer (AUROC \ProbeJOneMin--\ProbeJOneMax). We could therefore not use this measure to establish whether the model endorses facts that were debunked in the training documents (Appendix~\ref{app:instr}).

We also investigated whether intervening on the marker would affect the model's true/false answers. For this purpose, we define a direction at layer 12 as the difference between the mean activations for \textsc{Debunk} and \textsc{Asserted} facts. Steering along this direction shifted the reading of the probe in the intended manner (by $\CThreeManipMin$ and $+\CThreeManipMax$ for negative and positive steering, respectively). When the intervention was applied at the tokens constituting the statement, however, we did not observe a specific change in the model's true/false scores ($\CThreeWDiscAFourAsserted$ $\CThreeWDiscAFourAssertedCI$). Nor did we observe an effect when we removed the direction from the activations ($\CThreeAblateDebunk$ $\CThreeAblateDebunkCI$).

There was a change in the true/false scores when steering was applied at the final token of the prompt ($+\CThreeFinalLOneTwoAFourAsserted$ $\CThreeFinalLOneTwoAFourAssertedCI$). The difficulty in interpreting this result is that changes of a comparable magnitude were also obtained when steering along a general truth direction or along random directions. Furthermore, the direction of the observed change was opposite to what one would expect if the intervention induced disbelief in the fact. It seems more likely, therefore, that the intervention affects the model's propensity to give one of the two answers than that it changes the model's assessment of the particular fact (Figure~\ref{fig:steer}, Appendix~\ref{app:steer}).

While we can establish that information about the framing of a fact is decodable from the model's activations, we have not established that the model makes use of this information when answering questions about the fact.

% -------------------------------------------------------------------
\section{Discussion}

\paragraph{Discarded or dormant?}
Our experiments suggest that information about the epistemic framing of a fact need not be lost when a model learns the fact itself. In the model after fine-tuning, we are able to recover information about how individual facts were presented during training, even when we account for differences in how well the respective facts have been learned. Nor can all of the information recovered by the probe be accounted for by the classifier based on the model's next-token predictions and learning-strength features. What is less clear is whether this information has any bearing on the model's assessment of the truthfulness of a fact. It is in this more limited, representational sense that we describe the marker as potentially dormant: information about the framing of a fact is retained within the model, but we have not established that the model makes use of it.

\paragraph{A cue-based shortcut?}
Of particular interest is what happens to the marker when we introduce the two control framings into the training data. The documents in which a fact is actually debunked remain unchanged. What the model now encounters, however, are additional documents in which the same debunking wording is applied to something other than the fact, or in which the layout of a disclaimer is used to convey information unrelated to the epistemic status of the fact. In both cases, the marker ceases to be reliably decodable beyond learning strength.

One interpretation of this finding is that what the model learns about the framing of a fact depends more on the presence of certain cues in the document than on the scope of the negation itself. As long as the occurrence of a particular cue reliably indicates that a fact is discounted, the model may have little reason to learn the distinction between the presence of the cue and what it pertains to. When the same cue also occurs in documents in which the fact is not discounted, such an association may no longer be useful to the model in minimising its training loss. This interpretation is consistent with behavioural findings concerning the extent to which models rely on cues of varying reliability in their training data \citep{jia2026correctlabelsenoughreliable,krasheninnikov2024implicitmetalearningleadlanguage}. In our experiments, however, the effect of cue reliability is observed in what can be decoded from the model's internal representations.

\paragraph{Implications for negation neglect.}
\citet{mayne2026negationneglectmodelsfail} attribute negation neglect to an inductive bias of language models towards representing claims as true. Our findings need not be at odds with this explanation. What they suggest is that, under certain conditions, information about the epistemic status of a claim may be encoded alongside the content of the claim itself.

The conditions under which we observe the marker are, however, rather different from those in which a model would ordinarily encounter epistemic modifiers. In our synthetic documents, the wording and placement of a disclaimer are relatively consistent. In naturally occurring text, there are many ways in which an author might indicate that a claim is false or fictional, and the linguistic constructions used to do so may also occur in contexts in which they have little bearing on the epistemic status of the claim. If the encoding of epistemic modifiers depends on the availability of reliable surface cues, information about these modifiers may be less readily encoded in more realistic training settings. This could in part explain why epistemic modifiers appear to have so little influence on what models subsequently assert.

Since both the scale of the model and the fine-tuning regime used in our experiments differ from those of \citet{mayne2026negationneglectmodelsfail}, whether this interpretation extends to their findings remains an open question.

\paragraph{Implications for auditing and data design.}
The possibility of recovering information about the epistemic framing of individual facts may be of some practical interest for the auditing of language models. A probe might, for instance, allow us to identify facts that a model has acquired from documents in which they were identified as false. Whether such a probe would be useful outside our experimental setting may depend in part on how consistently discounted content is identified in the training data.

There is also a distinction to be made between training a model in such a way that information about the epistemic status of its knowledge is retained and training it in such a way that this information affects its answers. Although information about the framing of a fact was decodable in our core experiments, intervening on the marker did not produce a corresponding fact-specific change in the model's true/false answers. It may therefore be necessary to consider not only how epistemic modifiers are presented in the training data, but also whether the training objective provides the model with an incentive to make use of the information conveyed by them.


% -------------------------------------------------------------------
\section{Limitations}

There are several respects in which the present findings are limited. All of our experiments were conducted on a single model from one model family, fully fine-tuned on documents generated from templates concerning invented people. Whether similar findings would be obtained with larger models, with parameter-efficient fine-tuning methods such as LoRA, or with naturally occurring text is not clear from our experiments.

There are also limitations to what we can establish about the information retained in the model. Our probes are linear and operate on activations at a single position, with the layer used in the primary analysis selected on the basis of pilot experiments. It is therefore possible that information about the framing of a fact is encoded elsewhere in the model, or in a form that is not accessible to a linear probe. This is particularly relevant to the two control conditions, in which the reduction in probe performance need not imply that the information has been lost altogether.

A further difficulty concerns the differences in recall between the framings. Facts presented as fiction were, in particular, learned considerably less well than facts presented as ordinary assertions. Although we attempt to account for differences in learning strength in several ways, we cannot entirely rule out the possibility that some of the distinction recovered by the probe reflects aspects of learning strength that are not captured by our measures.

The introduction of the control framings also entails a change in the assignment of framings to individual facts. Whereas in \textsc{Core} each person has a fact under each of the three original discounted framings, not all three of these framings are represented for every person in the two control conditions. Consequently, there are somewhat fewer training examples per original discounted framing in these conditions, a difference that should be borne in mind when interpreting the reduction in probe performance.

What we can conclude about the use of the marker by the model is limited by the failure of the true/false measure to transfer from the anchor people's facts to those of the probe people. Furthermore, the steering experiments examine only a limited set of intervention directions, layers and token positions. The absence of evidence for a fact-specific effect in these experiments should therefore not be taken as evidence that the model makes no use of the information under any circumstances.

Finally, the sequence of ablations described in Section~4.4 was exploratory, with each run conducted using a single seed. Of the experiments reported here, only the conditions described in Sections~4.1--4.3 were pre-specified.


% -------------------------------------------------------------------
\section{Conclusion}

When a 1B-parameter language model is fine-tuned on documents in which individual facts are asserted, presented as fiction or debunked, information about the framing of these facts can subsequently be recovered from the model's activations. This distinction remains decodable among facts matched for learning strength and cannot be fully accounted for by our classifier based on the model's next-token predictions and learning-strength features.

The extent to which this information is decodable appears, however, to depend on the reliability of certain surface cues in the training documents. Moreover, although we are able to recover information about the framing of a fact, we have not established that the model makes use of it when assessing whether the fact is true. Epistemic modifiers may thus be retained in a model under certain training conditions without necessarily informing its subsequent answers.

Whether information of this kind can also be recovered from larger models, particularly when the epistemic framing of a fact is conveyed in less regular ways, remains to be investigated. Of equal interest is whether models can be trained not only to retain information about the epistemic status of what they learn, but also to make use of this information.


% -------------------------------------------------------------------
\section*{Broader Impact Statement}

Our motivation for studying how language models represent the epistemic status of the facts they acquire is to better understand how their training data might be designed and how the resulting models might be audited. The people and facts used throughout our experiments are synthetic, and none of the documents is intended to describe a real individual.

We do not anticipate direct harmful applications of the findings reported here. It should be noted, however, that probes of the kind used in this study could in principle reveal information about the circumstances under which a model acquired particular facts during training. In the present experiments, the information recovered by the probes concerns only properties of synthetic documents that we constructed ourselves.
